\honest{The Tragedy of Group Selectionism}{Eliezer Yudkowsky}{2007}

Before 1966, serious biologists held hypotheses ``that we would now regard as magical thinking.'' \nb{1966 is the year of George C.~Williams's \textsc{Adaptation and Natural Selection}, which with papers by Maynard Smith and Perrins cast serious doubt on group selection. The post names none of them.} Wynne-Edwards, Allee and Brereton believed that predators ``would voluntarily restrain their breeding'' so as not to exhaust their prey. \nb{In Wynne-Edwards's own summary, the restraint is imposed by territories and peck-orders: ``Individuals are made to compete for the right to feed,'' and ``surplus adults are inhibited from breeding.''}

A gene for restraint cannot spread by helping others; its effect must cause more copies of itself. Our sense of aesthetics cries that something should have been done when foxes eat all the rabbits and starve, as a human building a toy ecology would add a breeding restrainer. The temptation is then to find an argument that Nature wants the same thing for its own reasons. Individual selection does not favour restraint. (An addendum: it can favour four big eggs over eight small ones, if that is the individual's best trade-off, but not to spare shared resources.)

But if the species lives in small, isolated groups, restrained groups might survive and recolonize the land of groups that crashed. This was ``possible but very difficult.'' The benefits of restraint go to restrained and unrestrained alike, since unrestrained foxes and their many cubs eat the rabbits left over. Restraint spreads only if its cost to the donor, divided by its benefit to others, is less than the average relatedness between a fox and the neighbours who benefit. A simulation I link shows how demanding this is: with a 3\% cost, pure altruist groups twice as fit as selfish ones, groups of 25 and 20\% of deaths replaced by migrants, altruism merely survives alongside selfishness. Double the group size, or double the cost, and selfishness wins. For a cost above 10\%, groups must have about five members, which is implausible for foxes.

The group selectionists lost. The deciding blow was observation, although ``I forget the exact species of dispute'': predators did not restrain their breeding, and predator-prey systems crash all the time. Group selection later revived in a very different form. Population structure does create some group selection pressure, which your mathematics must include, though it need not beat individual selection. And evolved enforcement mechanisms change the game.

Then Michael Wade selected laboratory populations of flour beetles for small size. Did they restrain their breeding? ``No; the adults adapted to cannibalize eggs and larvae, especially female larvae.'' \nb{Wade's 1976 abstract lists four traits that changed: fecundity, developmental time, body weight and cannibalism. It also reports a fast, large response to group selection, and the experiment is cited, for example in the Stanford Encyclopedia of Philosophy, as evidence that group selection can be effective. The debate over group and multilevel selection continues.} Selecting for small populations would, of course, favour eating other individuals' children, not restraining one's own breeding. This is ``massively obvious in retrospect.'' \nb{In ``Hindsight Devalues Science,'' Yudkowsky had warned that hindsight makes findings seem less surprising than they were. There was better evidence the post does not use: cannibalism was known to keep flour beetle numbers down before Wade's work.}

My diagnosis: a missed third alternative, from a conclusion decided in advance, a fake justification and motivated stopping. The group selectionists started from ``the beautiful idea'' of nature in harmony, searched for a mechanism, and, knowing what they wanted it to produce, never asked neutrally what it would produce. Had they done so, they would have thought of cannibalism first. \nb{The post gives no evidence of how any of them reasoned. Wynne-Edwards's own account says the idea came while he was reviewing a book on animal numbers, and his picture included surplus animals that ``emigrate, or get killed by predators, disease, or starvation.''}

The lesson, as in ``Einstein's Arrogance,'' is that in a large space of answers nearly all the work goes into singling one out, so an answer singled out for its beauty may be doomed from the start. In principle you could still weigh the evidence neutrally and un-believe. In practice their conclusion came from aesthetics and Nature's from selection, two processes with no reason to agree. Nature does not yield to persuasive arguments; as J. R. Molloy said, it ``absolutely refuses to yield to the most persuasive rationalizations of humans.'' That is why I recommend evolutionary biology as training against rationalization. Engineers need no training against thinking electrons have minds, but natural selection produces purposes alien to ours, and students are warned about it. That is good practice for thinking about other mind-like processes that do not work as we do.
